\documentclass{subfiles}
\begin{document}
    In the previous section, we've defined a code as a map from the source alphabet 
        to a new one. The question is: are all code equal? 
        It's easely understood that this is not the case; 
        some code are in fact useless to our needs.
        For instance, if we define a code that maps two distinct character to the same 
        codeword, the receiver will have a hard time decoding the text.
        A code whose codewords are all distinct is called \emph{non-singular}. 
\end{document}
